\documentclass[10pt,a4paper,onecolumn]{article}
\usepackage{marginnote}
\usepackage{graphicx}
\usepackage{xcolor}
\usepackage{authblk,etoolbox}
\usepackage{titlesec}
\usepackage{calc}
\usepackage{tikz}
\usepackage{hyperref}
\hypersetup{colorlinks,breaklinks,
            urlcolor=[rgb]{0.0, 0.5, 1.0},
            linkcolor=[rgb]{0.0, 0.5, 1.0}}
\usepackage{caption}
\usepackage{tcolorbox}
\usepackage{amssymb,amsmath}
\usepackage{ifxetex,ifluatex}
\usepackage{seqsplit}
\usepackage{xstring}

\usepackage{float}
\let\origfigure\figure
\let\endorigfigure\endfigure
\renewenvironment{figure}[1][2] {
    \expandafter\origfigure\expandafter[H]
} {
    \endorigfigure
}

\usepackage{fixltx2e} % provides \textsubscript
\usepackage[
  backend=biber,
%  style=alphabetic,
%  citestyle=numeric
]{biblatex}
\bibliography{paper.bib}

% --- Splitting \texttt --------------------------------------------------

\let\textttOrig=\texttt
\def\texttt#1{\expandafter\textttOrig{\seqsplit{#1}}}
\renewcommand{\seqinsert}{\ifmmode
  \allowbreak
  \else\penalty6000\hspace{0pt plus 0.02em}\fi}


% --- Pandoc does not distinguish between links like [foo](bar) and
% --- [foo](foo) -- a simplistic Markdown model.  However, this is
% --- wrong:  in links like [foo](foo) the text is the url, and must
% --- be split correspondingly.
% --- Here we detect links \href{foo}{foo}, and also links starting
% --- with https://doi.org, and use path-like splitting (but not
% --- escaping!) with these links.
% --- Another vile thing pandoc does is the different escaping of
% --- foo and bar.  This may confound our detection.
% --- This problem we do not try to solve at present, with the exception
% --- of doi-like urls, which we detect correctly.


\makeatletter
\let\href@Orig=\href
\def\href@Urllike#1#2{\href@Orig{#1}{\begingroup
    \def\Url@String{#2}\Url@FormatString
    \endgroup}}
\def\href@Notdoi#1#2{\def\tempa{#1}\def\tempb{#2}%
  \ifx\tempa\tempb\relax\href@Urllike{#1}{#2}\else
  \href@Orig{#1}{#2}\fi}
\def\href#1#2{%
  \IfBeginWith{#1}{https://doi.org}%
  {\href@Urllike{#1}{#2}}{\href@Notdoi{#1}{#2}}}
\makeatother

% definitions for citeproc citations
\NewDocumentCommand\citeproctext{}{}
\NewDocumentCommand\citeproc{mm}{%
  \begingroup\def\citeproctext{#2}\cite{#1}\endgroup}
\makeatletter
 % allow citations to break across lines
 \let\@cite@ofmt\@firstofone
 % avoid brackets around text for \cite:
 \def\@biblabel#1{}
 \def\@cite#1#2{{#1\if@tempswa , #2\fi}}
\makeatother
\newlength{\cslhangindent}
\setlength{\cslhangindent}{1.5em}
\newlength{\csllabelwidth}
\setlength{\csllabelwidth}{3em}
\newenvironment{CSLReferences}[2] % #1 hanging-indent, #2 entry-spacing
 {\begin{list}{}{%
  \setlength{\itemindent}{0pt}
  \setlength{\leftmargin}{0pt}
  \setlength{\parsep}{0pt}
  % turn on hanging indent if param 1 is 1
  \ifodd #1
   \setlength{\leftmargin}{\cslhangindent}
   \setlength{\itemindent}{-1\cslhangindent}
  \fi
  % set entry spacing
  \setlength{\itemsep}{#2\baselineskip}}}
 {\end{list}}
\usepackage{calc}
\newcommand{\CSLBlock}[1]{\hfill\break\parbox[t]{\linewidth}{\strut\ignorespaces#1\strut}}
\newcommand{\CSLLeftMargin}[1]{\parbox[t]{\csllabelwidth}{\strut#1\strut}}
\newcommand{\CSLRightInline}[1]{\parbox[t]{\linewidth - \csllabelwidth}{\strut#1\strut}}
\newcommand{\CSLIndent}[1]{\hspace{\cslhangindent}#1}

% --- Page layout -------------------------------------------------------------
\usepackage[top=3.5cm, bottom=3cm, right=1.5cm, left=1.0cm,
            headheight=2.2cm, reversemp, includemp, marginparwidth=4.5cm]{geometry}

% --- Default font ------------------------------------------------------------
% \renewcommand\familydefault{\sfdefault}

% --- Style -------------------------------------------------------------------
\renewcommand{\bibfont}{\small \sffamily}
\renewcommand{\captionfont}{\small\sffamily}
\renewcommand{\captionlabelfont}{\bfseries}

% --- Section/SubSection/SubSubSection ----------------------------------------
\titleformat{\section}
  {\normalfont\sffamily\Large\bfseries}
  {}{0pt}{}
\titleformat{\subsection}
  {\normalfont\sffamily\large\bfseries}
  {}{0pt}{}
\titleformat{\subsubsection}
  {\normalfont\sffamily\bfseries}
  {}{0pt}{}
\titleformat*{\paragraph}
  {\sffamily\normalsize}


% --- Header / Footer ---------------------------------------------------------
\usepackage{fancyhdr}
\pagestyle{fancy}
\fancyhf{}
%\renewcommand{\headrulewidth}{0.50pt}
\renewcommand{\headrulewidth}{0pt}
\fancyhead[L]{\hspace{-0.75cm}\includegraphics[width=5.5cm]{.joss/logo.png}}
\fancyhead[C]{}
\fancyhead[R]{}
\renewcommand{\footrulewidth}{0.25pt}

\fancyfoot[L]{\parbox[t]{0.98\headwidth}{\footnotesize{\sffamily Kim et al., (2026). palsr:
Projected Actor Locations for spatial modeling of dyadic interactions
between moving
actors. \textit{Journal of Open Source Software}, VV(II), NNNN. \url{https://doi.org/XX.XXXXX/joss.0XXXX}}}}


\fancyfoot[R]{\sffamily \thepage}
\makeatletter
\let\ps@plain\ps@fancy
\fancyheadoffset[L]{4.5cm}
\fancyfootoffset[L]{4.5cm}

% --- Macros ---------

\definecolor{linky}{rgb}{0.0, 0.5, 1.0}

\newtcolorbox{repobox}
   {colback=red, colframe=red!75!black,
     boxrule=0.5pt, arc=2pt, left=6pt, right=6pt, top=3pt, bottom=3pt}

\newcommand{\ExternalLink}{%
   \tikz[x=1.2ex, y=1.2ex, baseline=-0.05ex]{%
       \begin{scope}[x=1ex, y=1ex]
           \clip (-0.1,-0.1)
               --++ (-0, 1.2)
               --++ (0.6, 0)
               --++ (0, -0.6)
               --++ (0.6, 0)
               --++ (0, -1);
           \path[draw,
               line width = 0.5,
               rounded corners=0.5]
               (0,0) rectangle (1,1);
       \end{scope}
       \path[draw, line width = 0.5] (0.5, 0.5)
           -- (1, 1);
       \path[draw, line width = 0.5] (0.6, 1)
           -- (1, 1) -- (1, 0.6);
       }
   }

% --- Title / Authors ---------------------------------------------------------
% patch \maketitle so that it doesn't center
\patchcmd{\@maketitle}{center}{flushleft}{}{}
\patchcmd{\@maketitle}{center}{flushleft}{}{}
% patch \maketitle so that the font size for the title is normal
\patchcmd{\@maketitle}{\LARGE}{\LARGE\sffamily}{}{}
% patch the patch by authblk so that the author block is flush left
\def\maketitle{{%
  \renewenvironment{tabular}[2][]
    {\begin{flushleft}}
    {\end{flushleft}}
  \AB@maketitle}}
\makeatletter
\renewcommand\AB@affilsepx{ \protect\Affilfont}
%\renewcommand\AB@affilnote[1]{{\bfseries #1}\hspace{2pt}}
\renewcommand\AB@affilnote[1]{{\bfseries #1}\hspace{3pt}}
\renewcommand{\affil}[2][]%
   {\newaffiltrue\let\AB@blk@and\AB@pand
      \if\relax#1\relax\def\AB@note{\AB@thenote}\else\def\AB@note{#1}%
        \setcounter{Maxaffil}{0}\fi
        \begingroup
        \let\href=\href@Orig
        \let\texttt=\textttOrig
        \let\protect\@unexpandable@protect
        \def\thanks{\protect\thanks}\def\footnote{\protect\footnote}%
        \@temptokena=\expandafter{\AB@authors}%
        {\def\\{\protect\\\protect\Affilfont}\xdef\AB@temp{#2}}%
         \xdef\AB@authors{\the\@temptokena\AB@las\AB@au@str
         \protect\\[\affilsep]\protect\Affilfont\AB@temp}%
         \gdef\AB@las{}\gdef\AB@au@str{}%
        {\def\\{, \ignorespaces}\xdef\AB@temp{#2}}%
        \@temptokena=\expandafter{\AB@affillist}%
        \xdef\AB@affillist{\the\@temptokena \AB@affilsep
          \AB@affilnote{\AB@note}\protect\Affilfont\AB@temp}%
      \endgroup
       \let\AB@affilsep\AB@affilsepx
}
\makeatother
\renewcommand\Authfont{\sffamily\bfseries}
\renewcommand\Affilfont{\sffamily\small\mdseries}
\setlength{\affilsep}{1em}


\ifnum 0\ifxetex 1\fi\ifluatex 1\fi=0 % if pdftex
  \usepackage[T1]{fontenc}
  \usepackage[utf8]{inputenc}

\else % if luatex or xelatex
  \ifxetex
    \usepackage{mathspec}
  \else
    \usepackage{fontspec}
  \fi
  \defaultfontfeatures{Ligatures=TeX,Scale=MatchLowercase}

\fi
% use upquote if available, for straight quotes in verbatim environments
\IfFileExists{upquote.sty}{\usepackage{upquote}}{}
% use microtype if available
\IfFileExists{microtype.sty}{%
\usepackage{microtype}
\UseMicrotypeSet[protrusion]{basicmath} % disable protrusion for tt fonts
}{}

\usepackage{hyperref}
\hypersetup{unicode=true,
            pdftitle={palsr: Projected Actor Locations for spatial modeling of dyadic interactions between moving actors},
            pdfborder={0 0 0},
            breaklinks=true}
\urlstyle{same}  % don't use monospace font for urls
\usepackage{color}
\usepackage{fancyvrb}
\newcommand{\VerbBar}{|}
\newcommand{\VERB}{\Verb[commandchars=\\\{\}]}
\DefineVerbatimEnvironment{Highlighting}{Verbatim}{commandchars=\\\{\}}
% Add ',fontsize=\small' for more characters per line
\newenvironment{Shaded}{}{}
\newcommand{\AlertTok}[1]{\textcolor[rgb]{1.00,0.00,0.00}{\textbf{#1}}}
\newcommand{\AnnotationTok}[1]{\textcolor[rgb]{0.38,0.63,0.69}{\textbf{\textit{#1}}}}
\newcommand{\AttributeTok}[1]{\textcolor[rgb]{0.49,0.56,0.16}{#1}}
\newcommand{\BaseNTok}[1]{\textcolor[rgb]{0.25,0.63,0.44}{#1}}
\newcommand{\BuiltInTok}[1]{\textcolor[rgb]{0.00,0.50,0.00}{#1}}
\newcommand{\CharTok}[1]{\textcolor[rgb]{0.25,0.44,0.63}{#1}}
\newcommand{\CommentTok}[1]{\textcolor[rgb]{0.38,0.63,0.69}{\textit{#1}}}
\newcommand{\CommentVarTok}[1]{\textcolor[rgb]{0.38,0.63,0.69}{\textbf{\textit{#1}}}}
\newcommand{\ConstantTok}[1]{\textcolor[rgb]{0.53,0.00,0.00}{#1}}
\newcommand{\ControlFlowTok}[1]{\textcolor[rgb]{0.00,0.44,0.13}{\textbf{#1}}}
\newcommand{\DataTypeTok}[1]{\textcolor[rgb]{0.56,0.13,0.00}{#1}}
\newcommand{\DecValTok}[1]{\textcolor[rgb]{0.25,0.63,0.44}{#1}}
\newcommand{\DocumentationTok}[1]{\textcolor[rgb]{0.73,0.13,0.13}{\textit{#1}}}
\newcommand{\ErrorTok}[1]{\textcolor[rgb]{1.00,0.00,0.00}{\textbf{#1}}}
\newcommand{\ExtensionTok}[1]{#1}
\newcommand{\FloatTok}[1]{\textcolor[rgb]{0.25,0.63,0.44}{#1}}
\newcommand{\FunctionTok}[1]{\textcolor[rgb]{0.02,0.16,0.49}{#1}}
\newcommand{\ImportTok}[1]{\textcolor[rgb]{0.00,0.50,0.00}{\textbf{#1}}}
\newcommand{\InformationTok}[1]{\textcolor[rgb]{0.38,0.63,0.69}{\textbf{\textit{#1}}}}
\newcommand{\KeywordTok}[1]{\textcolor[rgb]{0.00,0.44,0.13}{\textbf{#1}}}
\newcommand{\NormalTok}[1]{#1}
\newcommand{\OperatorTok}[1]{\textcolor[rgb]{0.40,0.40,0.40}{#1}}
\newcommand{\OtherTok}[1]{\textcolor[rgb]{0.00,0.44,0.13}{#1}}
\newcommand{\PreprocessorTok}[1]{\textcolor[rgb]{0.74,0.48,0.00}{#1}}
\newcommand{\RegionMarkerTok}[1]{#1}
\newcommand{\SpecialCharTok}[1]{\textcolor[rgb]{0.25,0.44,0.63}{#1}}
\newcommand{\SpecialStringTok}[1]{\textcolor[rgb]{0.73,0.40,0.53}{#1}}
\newcommand{\StringTok}[1]{\textcolor[rgb]{0.25,0.44,0.63}{#1}}
\newcommand{\VariableTok}[1]{\textcolor[rgb]{0.10,0.09,0.49}{#1}}
\newcommand{\VerbatimStringTok}[1]{\textcolor[rgb]{0.25,0.44,0.63}{#1}}
\newcommand{\WarningTok}[1]{\textcolor[rgb]{0.38,0.63,0.69}{\textbf{\textit{#1}}}}

% --- We redefined \texttt, but in sections and captions we want the
% --- old definition
\let\addcontentslineOrig=\addcontentsline
\def\addcontentsline#1#2#3{\bgroup
  \let\texttt=\textttOrig\addcontentslineOrig{#1}{#2}{#3}\egroup}
\let\markbothOrig\markboth
\def\markboth#1#2{\bgroup
  \let\texttt=\textttOrig\markbothOrig{#1}{#2}\egroup}
\let\markrightOrig\markright
\def\markright#1{\bgroup
  \let\texttt=\textttOrig\markrightOrig{#1}\egroup}


\usepackage{graphicx,grffile}
\makeatletter
\def\maxwidth{\ifdim\Gin@nat@width>\linewidth\linewidth\else\Gin@nat@width\fi}
\def\maxheight{\ifdim\Gin@nat@height>\textheight\textheight\else\Gin@nat@height\fi}
\makeatother
% Scale images if necessary, so that they will not overflow the page
% margins by default, and it is still possible to overwrite the defaults
% using explicit options in \includegraphics[width, height, ...]{}
\setkeys{Gin}{width=\maxwidth,height=\maxheight,keepaspectratio}
\IfFileExists{parskip.sty}{%
\usepackage{parskip}
}{% else
\setlength{\parindent}{0pt}
\setlength{\parskip}{6pt plus 2pt minus 1pt}
}
\setlength{\emergencystretch}{3em}  % prevent overfull lines
\providecommand{\tightlist}{%
  \setlength{\itemsep}{0pt}\setlength{\parskip}{0pt}}
\setcounter{secnumdepth}{0}
% Redefines (sub)paragraphs to behave more like sections
\ifx\paragraph\undefined\else
\let\oldparagraph\paragraph
\renewcommand{\paragraph}[1]{\oldparagraph{#1}\mbox{}}
\fi
\ifx\subparagraph\undefined\else
\let\oldsubparagraph\subparagraph
\renewcommand{\subparagraph}[1]{\oldsubparagraph{#1}\mbox{}}
\fi

\title{palsr: Projected Actor Locations for spatial modeling of dyadic
interactions between moving actors}

        \author[1]{Sangyeon Kim}
          \author[2]{Howard Liu}
          \author[3]{Bruce A. Desmarais}
    
      \affil[1]{Division of Communication and Media, Ewha Womans University, Republic of Korea}
      \affil[2]{Department of Political Science, University of South Carolina, USA}
      \affil[3]{Department of Political Science, Pennsylvania State University, USA}
  \date{\vspace{-5ex}}

\begin{document}
\maketitle

\marginpar{
  \sffamily\small

  {\bfseries DOI:} \href{https://doi.org/XX.XXXXX/joss.0XXXX}{\color{linky}{XX.XXXXX/joss.0XXXX}}

  \vspace{2mm}

  {\bfseries Software}
  \begin{itemize}
    \setlength\itemsep{0em}
    \item \href{https://github.com/openjournals/joss-reviews/issues/XXXX}{\color{linky}{Review}} \ExternalLink
    \item \href{https://github.com/bdesmarais/palsr}{\color{linky}{Repository}} \ExternalLink
    \item \href{https://doi.org/XX.XXXXX/zenodo.XXXXXXX}{\color{linky}{Archive}} \ExternalLink
  \end{itemize}

  \vspace{2mm}

  {\bfseries Submitted:} 13 June 2026\\
  {\bfseries Published:} 13 June 2026

  \vspace{2mm}
  {\bfseries License}\\
  Authors of papers retain copyright and release the work under a Creative Commons Attribution 4.0 International License (\href{https://creativecommons.org/licenses/by/4.0/}{\color{linky}{CC BY 4.0}}).
}

\section{Summary}\label{summary}

Many actors studied in the social sciences have no fixed geographic
location. Armed groups roam across a theater of conflict, firms
relocate, diplomats and international organizations shift their
operational focus, and migrating populations move in response to events.
For such \emph{moving actors}, the location at which two of them
interact is itself a meaningful, modelable outcome, but standard
spatial-interaction tools assume actors sit at fixed points and so
cannot be applied directly. The Projected Actor Location (PALS) method
(Kim et al. 2023) estimates where a mobile actor effectively ``is'' at any
point in time from the spatiotemporal history of its past interactions,
then uses these projected locations to model both the distance between
pairs of actors and the probability that they interact.

\texttt{palsr} is an R package (R Core Team 2024) that implements the
PALS workflow: constructing and validating dyadic event data, estimating
the smoothing parameters by minimizing great-circle prediction error,
projecting actor locations at arbitrary times, building the dyadic distance
covariates used in downstream interaction models, and quantifying
uncertainty through a nonparametric bootstrap with multiple-imputation
(Rubin's Rules) pooling. The performance-critical kernels (the Haversine
distance and the projection smoother) are implemented in C++ through
\texttt{Rcpp} (Eddelbuettel 2013), so that parameter estimation, which
repeatedly re-projects every actor across a grid of times, runs quickly
even on data sets with thousands of events.

\section{Statement of need}\label{statement-of-need}

The empirical study of geopolitical interaction (alliances, conflict,
cooperation, trade) is overwhelmingly \emph{dyadic}: the unit of analysis
is a pair of actors, and a central predictor of whether and how two actors
interact is the distance between them (Gleditsch and Ward 2001; Ward et al.
2007). When actors are states with fixed capitals or centroids, this
distance is trivial to compute. But many substantively important actors are
mobile: rebel groups, militias, terrorist organizations, peacekeeping
deployments, and non-governmental organizations. The fine-grained, geocoded event data that
now document them, such as the Armed Conflict Location and Event Data
project (Raleigh et al. 2010), record \emph{where each interaction
happened} but assign actors no persistent coordinates. Treating these
actors as fixed-location entities discards exactly the spatial dynamics
that drive their interactions.

Existing R spatial software does not do this. General geospatial toolkits
such as \texttt{sf} (Pebesma 2018) and spatial point-process libraries such
as \texttt{spatstat} (Baddeley and Turner 2005) operate on objects with
given coordinates; they provide no way to \emph{infer} a moving actor's
effective location,
let alone one tuned to predict future interactions. Conversely, dyadic
interaction models in political science and economics typically take
inter-actor distance as an exogenous input. Prior to PALS there was thus
no general, estimable method for assigning time-varying locations to
mobile actors in a way that is optimized for predicting their
interactions.

PALS was built for this setting, and the original article (Kim et al. 2023)
shows that PALS-based distances substantially improve the prediction of
subnational conflict relative to naive location measures. That contribution,
however, came only with replication scripts specific to a single
application. \texttt{palsr} packages the method as documented, tested
software that others can reuse. Applied researchers in political science,
conflict studies, economics, and geography can apply PALS to their own
dyadic event data without re-deriving the estimator; the package propagates
estimation uncertainty into downstream models through an event bootstrap,
and it ships both the original ACLED Nigeria data and a simulated-data generator,
so the full workflow can be learned and reproduced out of the box. To our knowledge it is the
first general-purpose software implementation of projected-location
modeling for moving actors.

\section{Method}\label{method}

For a focal actor \(i\) at prediction time \(t\), PALS forms a
recency-weighted mean of the locations of \(i\)'s own past events (the
\emph{focal} component) and a recency-weighted mean of the locations of
events involving \(i\)'s past interaction partners, or \emph{alters}
(the \emph{alter} component). The projected location is a convex
combination of the two,

\[
g_i(t) = (1-\pi)\sum_e W_i(e)\, g(e) \;+\; \pi \sum_e W_k(e)\, g(e),
\qquad \pi = \mathrm{logistic}(\gamma + \eta\, v),
\]

where event weights decay with age as in exponential smoothing (Gardner Jr.
2006), governed by \(\alpha\) for the focal actor and \(\beta\) for the
alters, and the mixing weight \(\pi\) depends through \(\gamma\) and
\(\eta\) on the focal actor's activity relative to its alters, summarized by
the event-count statistic \(v\). The four parameters are estimated by
``marching forward'' through time: every event is predicted using only
events strictly preceding it, and parameters are chosen to minimize the mean
Haversine distance between predicted and observed interaction locations. The
package also provides a one-parameter variant that fixes \(\pi = 0\) and
estimates only \(\alpha\), which is fast and frequently competitive.

Because the projected location is recomputed as time advances, each
actor traces a \emph{trajectory} through space rather than occupying a
fixed point. \autoref{fig:traj} illustrates this for four actors in the
bundled simulated data: PALS projects each actor's location at yearly
intervals, and the resulting paths drift through the cloud of observed
events.

\begin{figure}
\centering
\includegraphics[width=0.85\linewidth,height=\textheight,keepaspectratio,alt={PALS-projected locations of four major actors in the nigeria\_acled ACLED data (Boko Haram, the Nigerian military and police, and the Fulani Ethnic Militia), fit with the one-parameter model and projected at yearly intervals from 2009 to 2016. Grey points are the observed conflict events; coloured paths trace each actor's projected trajectory (arrows point forward in time). }]{figure-trajectories.png}
\caption{PALS-projected locations of four major actors in the
\protect\texttt{nigeria\_acled} ACLED data (Boko Haram, the Nigerian military and
police, and the Fulani Ethnic Militia), fit with the one-parameter model and
projected at yearly intervals from 2009 to 2016. Grey points are the observed
conflict events; coloured paths trace each actor's projected trajectory (arrows
point forward in time). \label{fig:traj}}
\end{figure}

\section{Key features}\label{key-features}

\begin{itemize}
\tightlist
\item
  \textbf{Validated data container} (\texttt{pal\_events()}) for dyadic,
  time-stamped, georeferenced events.
\item
  \textbf{Parameter estimation} (\texttt{estimate\_pals()}) for both the
  full four-parameter and the reduced one-parameter models, via
  Haversine-error minimization with a compiled objective.
\item
  \textbf{Projection} (\texttt{project\_pal()},
  \texttt{project\_pals()}) of actor locations at any set of times.
\item
  \textbf{Prediction and covariate construction}
  (\texttt{predict\_event\_locations()}, \texttt{pal\_distance()}) for
  interaction locations and dyadic distances, including a log transform
  for conflict-style specifications.
\item
  \textbf{Uncertainty quantification} (\texttt{bootstrap\_pals()},
  \texttt{pool\_rubin()}) via a nonparametric event bootstrap and
  Rubin's Rules (Rubin 1987).
\item
  \textbf{Reproducible example data} (\texttt{nigeria\_acled},
  \texttt{simulate\_conflict\_events()}), a vignette, and a full
  \texttt{testthat} suite that also checks the C++ kernels against
  pure-R reference implementations.
\end{itemize}

\section{Example}\label{example}

\begin{Shaded}
\begin{Highlighting}[]
\FunctionTok{library}\NormalTok{(palsr)}
\FunctionTok{data}\NormalTok{(nigeria\_acled)   }\CommentTok{\# 1,549 ACLED conflict events (Nigeria)}

\NormalTok{fit }\OtherTok{\textless{}{-}} \FunctionTok{estimate\_pals}\NormalTok{(nigeria\_acled, }\AttributeTok{model =} \StringTok{"one"}\NormalTok{)   }\CommentTok{\# estimate alpha}
\FunctionTok{coef}\NormalTok{(fit)}

\CommentTok{\# Project where each actor is on a given date.}
\FunctionTok{project\_pals}\NormalTok{(nigeria\_acled,}
\NormalTok{             }\AttributeTok{predict\_time =} \FunctionTok{as.Date}\NormalTok{(}\StringTok{"2015{-}12{-}01"}\NormalTok{),}
\NormalTok{             }\AttributeTok{params =}\NormalTok{ fit)}

\CommentTok{\# Build the dyadic distance covariate for an interaction model.}
\NormalTok{dyads }\OtherTok{\textless{}{-}} \FunctionTok{data.frame}\NormalTok{(}
\NormalTok{  }\AttributeTok{actor1 =} \StringTok{"Boko Haram"}\NormalTok{,}
\NormalTok{  }\AttributeTok{actor2 =} \StringTok{"Fulani Ethnic Militia (Nigeria)"}\NormalTok{,}
\NormalTok{  }\AttributeTok{time   =} \FunctionTok{as.Date}\NormalTok{(}\StringTok{"2014{-}06{-}01"}\NormalTok{)}
\NormalTok{)}
\FunctionTok{pal\_distance}\NormalTok{(nigeria\_acled, dyads, fit, }\AttributeTok{transform =} \StringTok{"log"}\NormalTok{)}

\CommentTok{\# Bootstrap uncertainty in the smoothing parameter.}
\FunctionTok{summary}\NormalTok{(}\FunctionTok{bootstrap\_pals}\NormalTok{(nigeria\_acled, }\AttributeTok{R =} \DecValTok{10}\NormalTok{, }\AttributeTok{model =} \StringTok{"one"}\NormalTok{, }\AttributeTok{seed =} \DecValTok{1}\NormalTok{))}
\end{Highlighting}
\end{Shaded}

\section{Availability and quality
control}\label{availability-and-quality-control}

\texttt{palsr} is developed openly on GitHub. The package ships with
documentation for every exported function, an introductory vignette that
walks through the full workflow on \texttt{nigeria\_acled}, and a
\texttt{testthat} suite covering data validation, estimation,
projection, distance construction, and uncertainty quantification. The
tests additionally verify the compiled C++ kernels against independent
pure-R reference implementations, and continuous integration runs
\texttt{R\ CMD\ check} across Linux, macOS, and Windows on every change.

\section{Acknowledgements}\label{acknowledgements}

This work was supported by the U.S. National Science Foundation under the
collaborative project ``Georelational Dynamics in Collective Action'' (awards
2446647 to B.A.D. and 2446646 to H.L.). We thank the editors and reviewers of
\emph{Political Science Research and Methods} for feedback on the method, and the
developers of \texttt{Rcpp} (Eddelbuettel 2013) for the tooling that underpins the
package's performance.

\section*{References}\label{references}
\addcontentsline{toc}{section}{References}

\protect\phantomsection\label{refs}
\begin{CSLReferences}{1}{1}
\bibitem[\citeproctext]{ref-spatstat}
Baddeley, Adrian, and Rolf Turner. 2005. {``Spatstat: An {R} Package for
Analyzing Spatial Point Patterns.''} \emph{Journal of Statistical
Software} 12 (6): 1--42. \url{https://doi.org/10.18637/jss.v012.i06}.

\bibitem[\citeproctext]{ref-rcpp}
Eddelbuettel, Dirk. 2013. \emph{Seamless {R} and {C++} Integration with
{Rcpp}}. Springer. \url{https://doi.org/10.1007/978-1-4614-6868-4}.

\bibitem[\citeproctext]{ref-gardner2006}
Gardner Jr., Everette S. 2006. {``Exponential Smoothing: The State of
the Art---Part {II}.''} \emph{International Journal of Forecasting} 22
(4): 637--66. \url{https://doi.org/10.1016/j.ijforecast.2006.03.005}.

\bibitem[\citeproctext]{ref-gleditsch2001}
Gleditsch, Kristian S., and Michael D. Ward. 2001. {``Measuring Space: A
Minimum-Distance Database and Applications to International Studies.''}
\emph{Journal of Peace Research} 38 (6): 739--58.
\url{https://doi.org/10.1177/0022343301038006006}.

\bibitem[\citeproctext]{ref-kim2023pals}
Kim, Sangyeon, Howard Liu, and Bruce A. Desmarais. 2023. {``Spatial
Modeling of Dyadic Geopolitical Interactions Between Moving Actors.''}
\emph{Political Science Research and Methods} 11 (3): 617--35.
\url{https://doi.org/10.1017/psrm.2022.6}.

\bibitem[\citeproctext]{ref-sf}
Pebesma, Edzer. 2018. {``Simple Features for {R}: Standardized Support
for Spatial Vector Data.''} \emph{The R Journal} 10 (1): 439--46.
\url{https://doi.org/10.32614/RJ-2018-009}.

\bibitem[\citeproctext]{ref-rcoreteam}
R Core Team. 2024. \emph{R: A Language and Environment for Statistical
Computing}. R Foundation for Statistical Computing.
\url{https://www.R-project.org/}.

\bibitem[\citeproctext]{ref-acled}
Raleigh, Clionadh, Andrew Linke, Håvard Hegre, and Joakim Karlsen. 2010.
{``Introducing {ACLED}: An Armed Conflict Location and Event Dataset.''}
\emph{Journal of Peace Research} 47 (5): 651--60.
\url{https://doi.org/10.1177/0022343310378914}.

\bibitem[\citeproctext]{ref-rubin1987}
Rubin, Donald B. 1987. \emph{Multiple Imputation for Nonresponse in
Surveys}. Wiley. \url{https://doi.org/10.1002/9780470316696}.

\bibitem[\citeproctext]{ref-ward2007}
Ward, Michael D., Randolph M. Siverson, and Xun Cao. 2007. {``Disputes,
Democracies, and Dependencies: A Reexamination of the Kantian Peace.''}
\emph{American Journal of Political Science} 51 (3): 583--601.
\url{https://doi.org/10.1111/j.1540-5907.2007.00269.x}.

\end{CSLReferences}

\end{document}
